% Extracted from LandEcon_Draft0.tex, lines 1248--1276.
% Insert in main document with: % Extracted from LandEcon_Draft0.tex, lines 1248--1276.
% Insert in main document with: % Extracted from LandEcon_Draft0.tex, lines 1248--1276.
% Insert in main document with: % Extracted from LandEcon_Draft0.tex, lines 1248--1276.
% Insert in main document with: \input{tables/tab_cross_study.tex}
% The tabular body is wrapped in \resizebox{\textwidth}{!}{...} to prevent margin overflow.
% Requires graphicx, which LandEcon_Draft0.tex already loads.

\begin{table}[h]
    \centering
    \caption{Cross-study summary of synthetic CV performance.}
    \resizebox{\textwidth}{!}{%
    \begin{threeparttable}
    \begin{tabular}{lcccc}
        \toprule
         & \multicolumn{2}{c}{Study 1 (NCES)} & \multicolumn{2}{c}{Study 2 (HWA)} \\
        \cmidrule(lr){2-3} \cmidrule(lr){4-5}
        Metric & Demo only & + Beliefs & Demo only & + Beliefs \\
        \midrule
        Best-LLM mean abs.\ deviation (pp)              & [X.XX] & [X.XX] & [X.XX] & [X.XX] \\
        Aggregate WTP error (\$)                        & [XXX]  & [XXX]  & [XXX]  & [XXX]  \\
        Best-LLM WTP error (\$)                         & [XXX]  & [XXX]  & [XXX]  & [XXX]  \\
        Synthetic-to-empirical bid coefficient ratio   & [0.XX] & [0.XX] & [0.XX] & [0.XX] \\
        Share of LLMs with WTP in original 95\% CI      & [X/10] & [X/10] & [X/10] & [X/10] \\
        \bottomrule
    \end{tabular}
    \begin{tablenotes}\footnotesize
    \item \textit{Notes:} Mean absolute deviation computed across bid
    levels and averaged across LLMs. Aggregate WTP refers to the
    equally-weighted ensemble across all ten models. The
    synthetic-to-empirical bid coefficient ratio is the cross-LLM mean of
    each model's bid coefficient divided by the corresponding empirical
    bid coefficient; values below one indicate flatter synthetic curves
    than empirical curves.
    \end{tablenotes}
    \end{threeparttable}
    }
\end{table}
% The tabular body is wrapped in \resizebox{\textwidth}{!}{...} to prevent margin overflow.
% Requires graphicx, which LandEcon_Draft0.tex already loads.

\begin{table}[h]
    \centering
    \caption{Cross-study summary of synthetic CV performance.}
    \resizebox{\textwidth}{!}{%
    \begin{threeparttable}
    \begin{tabular}{lcccc}
        \toprule
         & \multicolumn{2}{c}{Study 1 (NCES)} & \multicolumn{2}{c}{Study 2 (HWA)} \\
        \cmidrule(lr){2-3} \cmidrule(lr){4-5}
        Metric & Demo only & + Beliefs & Demo only & + Beliefs \\
        \midrule
        Best-LLM mean abs.\ deviation (pp)              & [X.XX] & [X.XX] & [X.XX] & [X.XX] \\
        Aggregate WTP error (\$)                        & [XXX]  & [XXX]  & [XXX]  & [XXX]  \\
        Best-LLM WTP error (\$)                         & [XXX]  & [XXX]  & [XXX]  & [XXX]  \\
        Synthetic-to-empirical bid coefficient ratio   & [0.XX] & [0.XX] & [0.XX] & [0.XX] \\
        Share of LLMs with WTP in original 95\% CI      & [X/10] & [X/10] & [X/10] & [X/10] \\
        \bottomrule
    \end{tabular}
    \begin{tablenotes}\footnotesize
    \item \textit{Notes:} Mean absolute deviation computed across bid
    levels and averaged across LLMs. Aggregate WTP refers to the
    equally-weighted ensemble across all ten models. The
    synthetic-to-empirical bid coefficient ratio is the cross-LLM mean of
    each model's bid coefficient divided by the corresponding empirical
    bid coefficient; values below one indicate flatter synthetic curves
    than empirical curves.
    \end{tablenotes}
    \end{threeparttable}
    }
\end{table}
% The tabular body is wrapped in \resizebox{\textwidth}{!}{...} to prevent margin overflow.
% Requires graphicx, which LandEcon_Draft0.tex already loads.

\begin{table}[h]
    \centering
    \caption{Cross-study summary of synthetic CV performance.}
    \resizebox{\textwidth}{!}{%
    \begin{threeparttable}
    \begin{tabular}{lcccc}
        \toprule
         & \multicolumn{2}{c}{Study 1 (NCES)} & \multicolumn{2}{c}{Study 2 (HWA)} \\
        \cmidrule(lr){2-3} \cmidrule(lr){4-5}
        Metric & Demo only & + Beliefs & Demo only & + Beliefs \\
        \midrule
        Best-LLM mean abs.\ deviation (pp)              & [X.XX] & [X.XX] & [X.XX] & [X.XX] \\
        Aggregate WTP error (\$)                        & [XXX]  & [XXX]  & [XXX]  & [XXX]  \\
        Best-LLM WTP error (\$)                         & [XXX]  & [XXX]  & [XXX]  & [XXX]  \\
        Synthetic-to-empirical bid coefficient ratio   & [0.XX] & [0.XX] & [0.XX] & [0.XX] \\
        Share of LLMs with WTP in original 95\% CI      & [X/10] & [X/10] & [X/10] & [X/10] \\
        \bottomrule
    \end{tabular}
    \begin{tablenotes}\footnotesize
    \item \textit{Notes:} Mean absolute deviation computed across bid
    levels and averaged across LLMs. Aggregate WTP refers to the
    equally-weighted ensemble across all ten models. The
    synthetic-to-empirical bid coefficient ratio is the cross-LLM mean of
    each model's bid coefficient divided by the corresponding empirical
    bid coefficient; values below one indicate flatter synthetic curves
    than empirical curves.
    \end{tablenotes}
    \end{threeparttable}
    }
\end{table}
% The tabular body is wrapped in \resizebox{\textwidth}{!}{...} to prevent margin overflow.
% Requires graphicx, which LandEcon_Draft0.tex already loads.

\begin{table}[h]
    \centering
    \caption{Cross-study summary of synthetic CV performance.}
    \resizebox{\textwidth}{!}{%
    \begin{threeparttable}
    \begin{tabular}{lcccc}
        \toprule
         & \multicolumn{2}{c}{Study 1 (NCES)} & \multicolumn{2}{c}{Study 2 (HWA)} \\
        \cmidrule(lr){2-3} \cmidrule(lr){4-5}
        Metric & Demo only & + Beliefs & Demo only & + Beliefs \\
        \midrule
        Best-LLM mean abs.\ deviation (pp)              & [X.XX] & [X.XX] & [X.XX] & [X.XX] \\
        Aggregate WTP error (\$)                        & [XXX]  & [XXX]  & [XXX]  & [XXX]  \\
        Best-LLM WTP error (\$)                         & [XXX]  & [XXX]  & [XXX]  & [XXX]  \\
        Synthetic-to-empirical bid coefficient ratio   & [0.XX] & [0.XX] & [0.XX] & [0.XX] \\
        Share of LLMs with WTP in original 95\% CI      & [X/10] & [X/10] & [X/10] & [X/10] \\
        \bottomrule
    \end{tabular}
    \begin{tablenotes}\footnotesize
    \item \textit{Notes:} Mean absolute deviation computed across bid
    levels and averaged across LLMs. Aggregate WTP refers to the
    equally-weighted ensemble across all ten models. The
    synthetic-to-empirical bid coefficient ratio is the cross-LLM mean of
    each model's bid coefficient divided by the corresponding empirical
    bid coefficient; values below one indicate flatter synthetic curves
    than empirical curves.
    \end{tablenotes}
    \end{threeparttable}
    }
\end{table}